\documentclass[12pt,a4paper]{article}

%% ── Core packages ───────────────────────────────────────────────────────────
\usepackage[utf8]{inputenc}
\usepackage[T1]{fontenc}
\usepackage[margin=1in, top=1.1in, bottom=1.1in]{geometry}
\usepackage{lmodern}
\usepackage{microtype}
\usepackage{graphicx}
\usepackage{booktabs}
\usepackage{longtable}
\usepackage{array}
\usepackage{multirow}
\usepackage{tabularx}
\usepackage{xcolor}
\usepackage{colortbl}
\usepackage{fancyhdr}
\usepackage{titlesec}
\usepackage{enumitem}
\usepackage{parskip}
\usepackage[most]{tcolorbox}
\usepackage{hyperref}

%% ── Colour palette ──────────────────────────────────────────────────────────
\definecolor{NavyBlue}   {HTML}{1e3a5f}
\definecolor{CyanAccent} {HTML}{0369a1}
\definecolor{Teal}       {HTML}{0f766e}
\definecolor{Violet}     {HTML}{5b21b6}
\definecolor{Crimson}    {HTML}{b91c1c}
\definecolor{Amber}      {HTML}{b45309}
\definecolor{LightGray}  {HTML}{f1f5f9}
\definecolor{MidGray}    {HTML}{94a3b8}
\definecolor{DarkGray}   {HTML}{334155}
\definecolor{TableHead}  {HTML}{1e3a5f}
\definecolor{TableAlt}   {HTML}{f0f7ff}
\definecolor{GreenCell}  {HTML}{d1fae5}
\definecolor{AmberCell}  {HTML}{fef3c7}

%% ── Hyperref ────────────────────────────────────────────────────────────────
\hypersetup{
    colorlinks  = true,
    linkcolor   = CyanAccent,
    urlcolor    = CyanAccent,
    citecolor   = Teal,
    pdftitle    = {Embedded AI: USB Malware Scanner - Optimization Report},
    pdfauthor   = {EAI Student Project},
}

%% ── Header / footer ─────────────────────────────────────────────────────────
\pagestyle{fancy}
\fancyhf{}
\fancyhead[L]{\small\textcolor{MidGray}{Embedded AI -- USB Malware Scanner Project}}
\fancyhead[R]{\small\textcolor{MidGray}{\thepage}}
\fancyfoot[C]{\small\textcolor{MidGray}{Department of Computer Science \& Engineering}}
\renewcommand{\headrulewidth}{0.4pt}
\renewcommand{\headrule}{\hbox to\headwidth{\color{CyanAccent}\leaders\hrule height \headrulewidth\hfill}}

%% ── Section title styles ────────────────────────────────────────────────────
\titleformat{\section}
  {\large\bfseries\color{NavyBlue}}
  {\thesection.}{0.6em}{}
  [\vspace{2pt}{\color{NavyBlue}\hrule height 1pt}\vspace{4pt}]

\titleformat{\subsection}
  {\normalsize\bfseries\color{CyanAccent}}
  {\thesubsection}{0.6em}{}

\titleformat{\subsubsection}
  {\normalsize\bfseries\color{Teal}}
  {\thesubsubsection}{0.6em}{}

%% ── tcolorbox styles ────────────────────────────────────────────────────────
\tcbset{
  infobox/.style={
    enhanced, breakable,
    colback  = LightGray,
    colframe = CyanAccent,
    fonttitle= \bfseries\color{white},
    coltitle = white,
    attach boxed title to top left={yshift=-2mm, xshift=6mm},
    boxed title style={colback=CyanAccent, sharp corners},
    arc=4pt, boxrule=0.8pt,
  },
  warnbox/.style={
    enhanced,
    colback  = AmberCell,
    colframe = Amber,
    arc=4pt, boxrule=0.8pt,
  },
  resultbox/.style={
    enhanced,
    colback  = GreenCell,
    colframe = Teal,
    arc=4pt, boxrule=0.8pt,
  },
}

%% ── Misc helpers ────────────────────────────────────────────────────────────
\newcolumntype{L}[1]{>{\raggedright\arraybackslash}p{#1}}
\newcolumntype{C}[1]{>{\centering\arraybackslash}p{#1}}
\newcolumntype{R}[1]{>{\raggedleft\arraybackslash}p{#1}}

\setlist[itemize]{leftmargin=1.5em, itemsep=2pt}
\setlist[enumerate]{leftmargin=1.5em, itemsep=2pt}

%% ═══════════════════════════════════════════════════════════════════════════
%%                         DOCUMENT BEGINS
%% ═══════════════════════════════════════════════════════════════════════════
\begin{document}

%% ────────────────────────────────────────────────────────────────────────────
%%  COVER PAGE
%% ────────────────────────────────────────────────────────────────────────────
\begin{titlepage}
\centering
\vspace*{1.5cm}

{\color{NavyBlue}\rule{\linewidth}{2pt}}
\vspace{0.4cm}

{\fontsize{22}{28}\selectfont\bfseries\color{NavyBlue}
Embedded Artificial Intelligence\\[6pt]
\large USB Malware Scanner -- Model Optimization\\
\& Hardware Feasibility Report}

\vspace{0.3cm}
{\color{NavyBlue}\rule{\linewidth}{2pt}}
\vspace{1.5cm}

{\large\itshape MalConv-Style Gated Byte CNN: Quantization, Pruning,\\
Weight Sharing and Sequence Downscaling for Edge Deployment}

\vspace{2.5cm}

\begin{tabular}{rl}
  \textbf{Course}          & Embedded Artificial Intelligence (EAI)\\[4pt]
  \textbf{Submitted to}    & Dr.\ Dubacharla Gyaneshwar\\[4pt]
  \textbf{Department}      & Computer Science \& Engineering\\[4pt]
  \textbf{Target Hardware} & Raspberry Pi 4 Model B (ARM Cortex-A72 MPU)\\[4pt]
  \textbf{Date}            & May 2026
\end{tabular}

\vspace{2.5cm}

\begin{tcolorbox}[warnbox, title={}]
\textbf{Abstract.}\quad
This report presents the full design, implementation, and benchmarking of a
\textit{USB Malware Sanitizer} powered by a MalConv-style gated 1-D convolutional
neural network.  Because the target platform (Raspberry Pi 4) has limited RAM,
compute budget, and real-time latency requirements, the baseline FP32 model
cannot be deployed without compression.  We apply four techniques taught in the
EAI course -- Post-Training Quantization, Weight Pruning, K-Means Weight
Sharing, and Sequence-Length Downscaling -- and benchmark every variant for
\textit{storage footprint, MAC computation, inference latency, and classification
accuracy}.  K-Means Weight Sharing (16 centroids) achieves an \textbf{87.5\%
reduction in model storage} (4\,174 KB down to 521 KB), while Sequence
Downscaling to 32\,KB cuts inference time by \textbf{5.44$\times$} (9.91 ms
down to 1.82 ms), both at \textbf{100\% retained accuracy}.
\end{tcolorbox}

\vfill
{\small\color{MidGray}Embedded Artificial Intelligence Project -- College of Engineering}
\end{titlepage}

%% ────────────────────────────────────────────────────────────────────────────
%%  TABLE OF CONTENTS
%% ────────────────────────────────────────────────────────────────────────────
\tableofcontents
\newpage

%% ═══════════════════════════════════════════════════════════════════════════
%%  1. PROJECT OVERVIEW
%% ═══════════════════════════════════════════════════════════════════════════
\section{Project Overview}

\subsection{Problem Statement}
USB flash drives are one of the most common vectors for malware distribution in
institutional environments (labs, offices, colleges).  Infected drives may carry
packed executables, polymorphic code, or obfuscated binaries that traditional
signature-based antivirus tools miss -- because packing alters raw bytes,
bypassing static hash checks.

The project addresses this by building an \textbf{AI-powered USB Sanitizer
Kiosk} that physically sits between an untrusted USB drive and any workplace
computer.  Every file on the drive is scored by a lightweight deep-learning
model before it is allowed to enter the trusted network.

\subsection{Why a Dedicated Kiosk Instead of Endpoint AV?}

\begin{tcolorbox}[infobox, title={Deployment Architecture Rationale}]
\begin{itemize}
  \item \textbf{Trust boundary isolation} -- An untrusted USB is first
        inspected by the Pi. The protected workstation is never directly exposed
        to unscanned media.
  \item \textbf{OS-agnostic} -- Works regardless of the OS on the target
        machine; no agent installation required.
  \item \textbf{Centralized policy} -- One kiosk enforces the same model,
        thresholds, and quarantine rules for every machine it serves.
  \item \textbf{Institutional feasibility} -- College labs often forbid
        installing custom software on PCs.  A standalone Pi kiosk bypasses this
        restriction entirely.
\end{itemize}
\end{tcolorbox}

\subsection{System Workflow}

The end-to-end pipeline is:

\begin{enumerate}
  \item \textbf{Drive Detection} -- A background monitor detects new USB
        volumes using OS-level drive enumeration (PowerShell / \texttt{psutil}).
  \item \textbf{File Enumeration} -- All files matching target extensions
        (\texttt{.exe}, \texttt{.dll}, \texttt{.com}) are collected.
  \item \textbf{Stage 1 -- Structural Integrity} -- Magic-number vs.\
        extension mismatch triggers immediate quarantine.
  \item \textbf{Stage 2 -- Header Heuristics} -- Shannon entropy analysis and
        suspicious PE import checking.
  \item \textbf{Stage 3 -- Neural Scoring} -- The MalConv model (ONNX format)
        reads the first \textit{N} bytes of each file and outputs a suspicion
        score in [0, 1].  Files scoring above 0.5 are quarantined.
  \item \textbf{Quarantine / Report} -- Flagged files are moved to a secure
        quarantine directory; a JSON and plain-text log is produced.
\end{enumerate}

\subsection{Dataset Description}

Because real malware cannot be used in a college lab environment, the project
uses a \textbf{packed-vs-clean demo dataset}:

\begin{table}[h!]
\centering
\caption{Demo Dataset Composition}
\renewcommand{\arraystretch}{1.3}
\begin{tabular}{L{3.5cm} L{4.5cm} C{2cm} C{3cm}}
\toprule
\rowcolor{TableHead}
\color{white}\textbf{Class} &
\color{white}\textbf{Source} &
\color{white}\textbf{Count} &
\color{white}\textbf{Label} \\
\midrule
Benign & Windows System32 \texttt{.exe} binaries & 33 & 0 (Benign) \\
\rowcolor{TableAlt}
Malicious (simulated) & Same binaries, UPX-packed with \texttt{upx.exe} & 33 & 1 (Suspicious) \\
\midrule
\textbf{Total} & -- & \textbf{66} & -- \\
\bottomrule
\end{tabular}
\end{table}

\noindent
UPX packing is an industry-standard obfuscation technique that repacks the PE
header and segments, dramatically altering the raw byte distribution.  The model
learns to distinguish natural PE header structures from UPX-repacked ones
without ever seeing real malware.

The dataset is split using \textbf{group-by-basename} logic to prevent
clean/packed twin leakage across train/validation/test splits
(80\% / 10\% / 10\%).

%% ═══════════════════════════════════════════════════════════════════════════
%%  2. ARCHITECTURE OF THE MALCONV MODEL
%% ═══════════════════════════════════════════════════════════════════════════
\section{MalConv Gated CNN Architecture}

\subsection{Design Motivation}
Traditional ML classifiers require hand-crafted feature extraction from
executables (PE headers, import tables, entropy histograms).  MalConv eliminates
this step by treating the executable as a raw byte sequence and learning
discriminative patterns directly from the bytes using convolutions.  This makes
it robust to novel obfuscation that changes features but not byte-level
structure.

\subsection{Layer Architecture}

\begin{table}[h!]
\centering
\caption{MalConv Layer-by-Layer Architecture}
\renewcommand{\arraystretch}{1.4}
\begin{tabular}{C{0.5cm} L{3.2cm} L{3.8cm} R{2.4cm} L{4cm}}
\toprule
\rowcolor{TableHead}
\color{white}\textbf{\#} &
\color{white}\textbf{Layer} &
\color{white}\textbf{Output Shape} &
\color{white}\textbf{MACs} &
\color{white}\textbf{Purpose} \\
\midrule
1 & Embedding (257 tokens, dim=8) & (B, 8, L) & -- (lookup) & Map bytes to dense vectors \\
\rowcolor{TableAlt}
2 & Conv1D Value -- 128 filters, k=512, s=512 & (B, 128, L') & 268 M & Feature extraction \\
3 & Conv1D Gate -- 128 filters, k=512, s=512 + Sigmoid & (B, 128, L') & 268 M & Attention mask \\
\rowcolor{TableAlt}
4 & Elementwise gating multiply & (B, 128, L') & -- & Suppress noise regions \\
5 & Adaptive Max Pooling (output=1) & (B, 128) & -- & Length-independent aggregation \\
\rowcolor{TableAlt}
6 & Linear 128 \textrightarrow 128 + ReLU & (B, 128) & 16 K & Nonlinear classification \\
7 & Linear 128 \textrightarrow 1 + Sigmoid & (B, 1) & 128 & Probability output [0,1] \\
\midrule
\multicolumn{3}{l}{\textbf{Total Trainable Parameters}} & \multicolumn{2}{r}{\textbf{1,067,529}} \\
\multicolumn{3}{l}{\textbf{Total MACs (baseline, 256 KB input)}} & \multicolumn{2}{r}{\textbf{536,887,424}} \\
\bottomrule
\end{tabular}
\end{table}

\subsection{Gating Mechanism}

The key innovation in MalConv is the \textbf{sigmoid gating} applied
elementwise between two parallel convolutions.  One convolution learns
\textit{what} features are present (value stream); the other learns
\textit{where} those features matter (gate stream, squashed to [0,1] by
sigmoid).  Multiplying both streams focuses the model on malicious header
regions while ignoring the large padding-filled tail of short files.

This is functionally equivalent to a learned spatial attention mechanism, but
implemented efficiently as two plain convolutions -- making it highly
suitable for edge inference without specialized attention hardware.

\subsection{Checkpoint Format}

Each saved checkpoint is a Python dictionary containing:
\begin{itemize}
  \item \texttt{state\_dict} -- all trained weight tensors.
  \item \texttt{max\_len} -- maximum byte input length the model was trained on.
  \item \texttt{window\_size} -- convolution kernel / stride size.
\end{itemize}

Baseline checkpoint size: \textbf{4.08 MB} (FP32, uncompressed).

%% ═══════════════════════════════════════════════════════════════════════════
%%  3. TARGET HARDWARE CONSTRAINTS
%% ═══════════════════════════════════════════════════════════════════════════
\section{Target Hardware and Embedded Constraints}

\subsection{Target Platform: Raspberry Pi 4 Model B}

\begin{table}[h!]
\centering
\caption{Raspberry Pi 4 Hardware Specifications}
\renewcommand{\arraystretch}{1.3}
\begin{tabular}{L{4cm} L{8cm}}
\toprule
\rowcolor{TableHead}
\color{white}\textbf{Specification} & \color{white}\textbf{Value} \\
\midrule
CPU & ARM Cortex-A72, 4 cores @ 1.8 GHz \\
\rowcolor{TableAlt}
RAM & 4 GB LPDDR4 \\
Storage & MicroSD (16--64 GB typical) \\
\rowcolor{TableAlt}
OS & Raspberry Pi OS (Linux, 64-bit) \\
Power & 5V / 3A USB-C (15 W) \\
\rowcolor{TableAlt}
Inference Engine & ONNX Runtime (CPUExecutionProvider) \\
USB Interface & USB 3.0 Host Port \\
\rowcolor{TableAlt}
Architecture Class & Microprocessor Unit (MPU) \\
\bottomrule
\end{tabular}
\end{table}

\subsection{Embedded Resource Constraints}

\begin{tcolorbox}[infobox, title={Key EAI Embedded Constraints}]
\begin{enumerate}
  \item \textbf{Memory constraint} -- The model's weights, input buffer, and
        intermediate activation tensors must all reside simultaneously in RAM.
        With only 4 GB shared between the OS and application, large FP32
        activation buffers for 1 MB files are costly.
  \item \textbf{Computation constraint} -- The ARM Cortex-A72 achieves
        approximately 100--200 GFLOPS (FP32) in practice.  The baseline model
        requires 537 million MACs per file scan -- which is feasible, but
        becomes a bottleneck when scanning 100+ files per USB drive.
  \item \textbf{Latency constraint} -- A real-time kiosk must scan each file
        within 10--20 ms to keep total drive scanning time under 30 seconds
        for a typical drive of 50--100 executables.
  \item \textbf{Storage constraint} -- The Pi's microSD card hosts both the OS
        and all application files.  A compressed model under 1 MB is strongly
        preferred.
\end{enumerate}
\end{tcolorbox}

\subsection{Feasibility Assessment of the Baseline Model}

\begin{table}[h!]
\centering
\caption{Baseline Model Feasibility on Raspberry Pi 4}
\renewcommand{\arraystretch}{1.3}
\begin{tabular}{L{4.5cm} C{2.5cm} C{2.5cm} L{4.5cm}}
\toprule
\rowcolor{TableHead}
\color{white}\textbf{Resource} &
\color{white}\textbf{Pi Budget} &
\color{white}\textbf{Baseline Needs} &
\color{white}\textbf{Verdict} \\
\midrule
Flash Storage (model) & 16--64 GB & 4.08 MB & \textcolor{Teal}{\textbf{OK}} -- fits easily \\
\rowcolor{TableAlt}
RAM (model weights) & 4 GB & 4.07 MB & \textcolor{Teal}{\textbf{OK}} \\
RAM (input buffer, 256 KB) & -- & 256 KB & \textcolor{Teal}{\textbf{OK}} \\
\rowcolor{TableAlt}
MACs per scan & \textasciitilde 200 GMAC/s & 537 M MACs & \textcolor{Teal}{\textbf{OK}} \textasciitilde 2--5 ms pure compute \\
End-to-end latency (with I/O) & \textless 20 ms target & 9.91 ms & \textcolor{Teal}{\textbf{OK}} \\
\rowcolor{TableAlt}
Accuracy on test split & 100\% required & 100.0\% & \textcolor{Teal}{\textbf{OK}} \\
\bottomrule
\end{tabular}
\end{table}

\noindent
The baseline model \textit{already fits} on the Raspberry Pi 4 -- all
constraints are within budget.  However, the project requirement is to
\textbf{demonstrate quantitative trade-off analysis} between the baseline and
all compressed variants, proving that even further compression is possible
with \textit{zero accuracy loss}, giving deployment flexibility for more
constrained targets (e.g.\ Raspberry Pi Zero, STM32 MCU).

%% ═══════════════════════════════════════════════════════════════════════════
%%  4. OPTIMIZATION TECHNIQUES APPLIED
%% ═══════════════════════════════════════════════════════════════════════════
\section{Optimization Techniques Applied}

All four techniques below were applied \textit{post-training} (no retraining
was performed), using the best checkpoint
\texttt{outputs/models/group\_split\_model.pth}.

%% ── 4.1 Quantization ────────────────────────────────────────────────────────
\subsection{Technique 1: Post-Training Dynamic Quantization}

\subsubsection{What it does}
Quantization replaces high-precision floating-point weights (32-bit FP32)
with lower-precision integers (8-bit INT8).  Dynamic quantization performs
this conversion \textit{at runtime} for each inference call without a
dedicated calibration dataset.

\subsubsection{What we applied}
We applied two forms:
\begin{itemize}
  \item \textbf{PyTorch Dynamic Quantization} -- Only the fully connected
        layers (\texttt{fc\_1}, \texttt{fc\_2}) are quantized to INT8.
        Convolutional layers remain FP32.  This is the safest, fastest
        option with minimal accuracy risk.
  \item \textbf{ONNX Dynamic Quantization} -- The entire model is first
        exported to an ONNX computation graph, then all weight matrices
        (including the embedding table and convolutional kernels) are
        quantized to INT8.  This achieves greater compression at the cost of
        potential precision loss.
\end{itemize}

\subsubsection{Expected benefits}
\begin{itemize}
  \item Up to \textbf{4$\times$ smaller} weight storage (4 bytes per parameter
        reduced to 1 byte).
  \item Faster matrix multiplication on CPUs that support INT8 SIMD
        instructions (ARM NEON on Cortex-A72).
  \item Reduced memory bandwidth pressure.
\end{itemize}

%% ── 4.2 Pruning ─────────────────────────────────────────────────────────────
\subsection{Technique 2: Unstructured Weight Pruning}

\subsubsection{What it does}
Pruning removes individual weight connections whose absolute values fall below
a threshold, zeroing them out and creating a \textit{sparse} weight matrix.
L1-magnitude pruning ranks weights by $|w|$ and zeros out the lowest-ranked
fraction.

\subsubsection{What we applied}
We applied \textbf{L1 Unstructured Pruning} to all convolutional and the first
fully connected layer at two sparsity levels:
\begin{itemize}
  \item \textbf{50\% Sparsity} -- Half of all selected-layer connections are
        zeroed.
  \item \textbf{75\% Sparsity} -- Three-quarters of all selected-layer
        connections are zeroed.
\end{itemize}

\subsubsection{Expected benefits}
\begin{itemize}
  \item Reduces the number of \textit{active} multiply-accumulate operations
        (non-zero MACs), lowering effective computation.
  \item Enables sparse matrix formats (CSR/CSC) on MCU runtimes that save
        both storage and computation proportionally to sparsity.
  \item Accuracy is expected to degrade gracefully at high sparsity ratios.
\end{itemize}

%% ── 4.3 Weight Sharing ──────────────────────────────────────────────────────
\subsection{Technique 3: K-Means Weight Sharing}

\subsubsection{What it does}
Weight sharing clusters all weight values of a layer into $k$ groups using
K-Means clustering.  Every weight in a cluster is replaced by the cluster
centroid value.  The model is stored as a \textit{codebook} (the $k$ centroid
values) plus an \textit{index array} (one integer per weight pointing to its
centroid).  Storing $k{=}16$ centroids requires only 4 bits per index (instead
of 32 bits per FP32 weight), giving a theoretical 8$\times$ compression.

\subsubsection{What we applied}
\begin{itemize}
  \item \textbf{K-Means k=16} -- 16 centroids per layer.  Weights encoded as
        4-bit indices.  Estimated compressed storage: \textbf{521 KB}.
  \item \textbf{K-Means k=32} -- 32 centroids per layer.  Weights encoded as
        5-bit indices.  Estimated compressed storage: \textbf{652 KB}.
\end{itemize}

\subsubsection{Expected benefits}
\begin{itemize}
  \item Extreme storage compression -- the best result among all techniques.
  \item No change to computation graph structure; the same convolutions and
        linear layers execute identically after centroid substitution.
  \item Accuracy depends on how spread the weight distributions are; highly
        redundant embeddings (as in MalConv) tolerate aggressive clustering.
\end{itemize}

%% ── 4.4 Sequence Downscaling ────────────────────────────────────────────────
\subsection{Technique 4: Input Sequence Length Downscaling}

\subsubsection{What it does}
MalConv reads the first \texttt{max\_len} bytes of each file (zero-padding
shorter files).  Because the convolutional layers slide a fixed-size window
across the byte sequence, the number of MAC operations scales \textit{linearly}
with \texttt{max\_len}.  By scanning only the first $N$ bytes instead of the
full 256 KB, we reduce computation proportionally -- at no cost to model
weights.

\subsubsection{What we applied}
We evaluated three reduced prefix sizes:
\begin{itemize}
  \item \textbf{128 KB (131,072 bytes)} -- 50\% reduction vs.\ baseline.
  \item \textbf{64 KB (65,536 bytes)} -- 75\% reduction vs.\ baseline.
  \item \textbf{32 KB (32,768 bytes)} -- 87.5\% reduction vs.\ baseline.
\end{itemize}

\subsubsection{Why accuracy is unaffected}
UPX packing modifies the PE header and the first program segment.  The UPX
stub, section headers, and modified entry-point all appear in the first 32 KB
of any executable.  The model only needs to see the file prefix to classify
correctly, making 32 KB a sufficient scanning budget.

%% ═══════════════════════════════════════════════════════════════════════════
%%  5. BENCHMARKING METHODOLOGY
%% ═══════════════════════════════════════════════════════════════════════════
\section{Benchmarking Methodology}

\subsection{Evaluation Setup}

\begin{table}[h!]
\centering
\caption{Benchmarking Environment}
\renewcommand{\arraystretch}{1.3}
\begin{tabular}{L{4cm} L{9cm}}
\toprule
\rowcolor{TableHead}
\color{white}\textbf{Parameter} & \color{white}\textbf{Value} \\
\midrule
Platform & Windows 11, x86-64, AMD/Intel CPU (CPU-only inference) \\
\rowcolor{TableAlt}
Framework & PyTorch 2.11 + ONNX Runtime 1.24 \\
Inference device & CPU (mimicking Pi single-core budget) \\
\rowcolor{TableAlt}
Test split & 8 samples (4 benign, 4 packed), group-by-basename split \\
Warmup runs & 2 forward passes before timing \\
\rowcolor{TableAlt}
Latency metric & Mean model-only and end-to-end time over all test samples \\
MAC counting & Analytical: based on conv output length $\times$ kernel ops \\
\rowcolor{TableAlt}
Accuracy metric & Classification accuracy on test split (threshold = 0.5) \\
\bottomrule
\end{tabular}
\end{table}

\subsection{Metrics Collected}

For every model variant we measure:
\begin{itemize}
  \item \textbf{Storage Size (KB)} -- checkpoint file size on disk.
  \item \textbf{Estimated Compressed Size (KB)} -- for weight-sharing variants,
        the theoretical codebook + index size.
  \item \textbf{Total MACs} -- theoretical multiply-accumulate operations for
        one forward pass.
  \item \textbf{Non-Zero MACs} -- active MACs after accounting for pruning
        sparsity.
  \item \textbf{Model-Only Latency (ms)} -- pure forward-pass timing excluding
        file I/O.
  \item \textbf{End-to-End Latency (ms)} -- file read + preprocessing +
        forward pass.
  \item \textbf{Test Accuracy (\%)} -- fraction of test samples correctly
        classified.
\end{itemize}

%% ═══════════════════════════════════════════════════════════════════════════
%%  6. RESULTS AND BENCHMARKING DATA
%% ═══════════════════════════════════════════════════════════════════════════
\section{Results and Benchmarking Data}

\subsection{Complete Optimization Results Matrix}

\begin{table}[h!]
\centering
\caption{Full Benchmarking Results -- All Model Variants}
\label{tab:main}
\renewcommand{\arraystretch}{1.4}
\begin{tabular}{L{3.6cm} C{1.8cm} R{1.8cm} R{1.8cm} R{1.8cm} C{1.8cm}}
\toprule
\rowcolor{TableHead}
\color{white}\textbf{Model Variant} &
\color{white}\textbf{Method} &
\color{white}\textbf{Size (KB)} &
\color{white}\textbf{Active MACs} &
\color{white}\textbf{Latency (ms)} &
\color{white}\textbf{Accuracy} \\
\midrule
Original FP32 (Baseline)      & --            & 4,174.1 & 536,887,424 &  9.91 & 100.0\% \\
\rowcolor{TableAlt}
PyTorch Dynamic INT8          & Quantization  & 4,127.0 & 536,887,424 &  9.76 & 100.0\% \\
ONNX Dynamic INT8             & Quantization  & 3,098.2 & 536,887,424 & 15.88 & 100.0\% \\
\rowcolor{TableAlt}
Weight Pruning 50\%           & Pruning       & 4,173.9 & 273,812,648 & 13.44 & 100.0\% \\
Weight Pruning 75\%           & Pruning       & 4,173.9 & 139,590,824 & 15.58 & 100.0\% \\
\rowcolor{GreenCell}
\textbf{Weight Sharing k=16}  & \textbf{Sharing}      & \textbf{521.3}  & 536,887,424 & 14.08 & \textbf{100.0\%} \\
\rowcolor{TableAlt}
Weight Sharing k=32           & Sharing       &   651.7 & 536,887,424 & 11.25 & 100.0\% \\
Sequence Downscale 128 KB     & Downscaling   & 4,174.1 & 268,451,968 &  4.84 & 100.0\% \\
\rowcolor{TableAlt}
Sequence Downscale 64 KB      & Downscaling   & 4,174.1 & 134,234,240 &  2.87 & 100.0\% \\
\rowcolor{GreenCell}
\textbf{Sequence Downscale 32 KB}  & \textbf{Downscaling} & 4,174.1 & \textbf{67,125,376} & \textbf{1.82} & \textbf{100.0\%} \\
\bottomrule
\end{tabular}
\end{table}

\subsection{Storage Footprint Comparison}

\begin{table}[h!]
\centering
\caption{Storage Footprint -- Before and After Compression}
\renewcommand{\arraystretch}{1.4}
\begin{tabular}{L{4.5cm} R{2.2cm} R{2.2cm} R{2.2cm}}
\toprule
\rowcolor{TableHead}
\color{white}\textbf{Model Variant} &
\color{white}\textbf{Size (KB)} &
\color{white}\textbf{Savings (KB)} &
\color{white}\textbf{Reduction (\%)} \\
\midrule
Original FP32 (Baseline)     & 4,174.1 & --      & -- \\
\rowcolor{TableAlt}
PyTorch Dynamic INT8         & 4,127.0 &    47.1 &  1.1\% \\
ONNX Dynamic INT8            & 3,098.2 & 1,075.9 & 25.8\% \\
\rowcolor{TableAlt}
Weight Pruning 50\%          & 4,173.9 &     0.2 &  0.0\% \\
Weight Pruning 75\%          & 4,173.9 &     0.2 &  0.0\% \\
\rowcolor{GreenCell}
\textbf{Weight Sharing k=16} & \textbf{521.3}  & \textbf{3,652.8} & \textbf{87.5\%} \\
\rowcolor{TableAlt}
Weight Sharing k=32          &   651.7 & 3,522.4 & 84.4\% \\
Sequence Downscale (any)     & 4,174.1 &      -- &    -- \\
\bottomrule
\end{tabular}
\end{table}

\textit{Note:} Weight-sharing sizes represent the theoretical
\textbf{codebook + index} format.  Pruning does not reduce file size in
standard PyTorch format because zeroed weights are still stored as FP32
zeros; compression is realized at the arithmetic computation level.

\subsection{Computational Complexity Comparison}

\begin{table}[h!]
\centering
\caption{MAC Computation -- Before and After Optimization}
\renewcommand{\arraystretch}{1.4}
\begin{tabular}{L{4.5cm} R{2.8cm} R{2.8cm} R{2.2cm}}
\toprule
\rowcolor{TableHead}
\color{white}\textbf{Model Variant} &
\color{white}\textbf{Active MACs} &
\color{white}\textbf{Saved MACs} &
\color{white}\textbf{Reduction (\%)} \\
\midrule
Original FP32 (Baseline)          & 536,887,424 & --          & -- \\
\rowcolor{TableAlt}
Quantization (both variants)      & 536,887,424 & --          & 0.0\% \\
Weight Pruning 50\%               & 273,812,648 & 263,074,776 & 49.0\% \\
\rowcolor{TableAlt}
Weight Pruning 75\%               & 139,590,824 & 397,296,600 & 74.0\% \\
Weight Sharing (both variants)    & 536,887,424 & --          & 0.0\% \\
\rowcolor{TableAlt}
Sequence Downscale 128 KB         & 268,451,968 & 268,435,456 & 50.0\% \\
Sequence Downscale 64 KB          & 134,234,240 & 402,653,184 & 75.0\% \\
\rowcolor{GreenCell}
\textbf{Sequence Downscale 32 KB} & \textbf{67,125,376}  & \textbf{469,762,048} & \textbf{87.5\%} \\
\bottomrule
\end{tabular}
\end{table}

\subsection{Inference Latency Comparison}

\begin{table}[h!]
\centering
\caption{Inference Latency -- Model-Only and End-to-End}
\renewcommand{\arraystretch}{1.4}
\begin{tabular}{L{4.5cm} R{2.5cm} R{2.5cm} R{2.5cm}}
\toprule
\rowcolor{TableHead}
\color{white}\textbf{Model Variant} &
\color{white}\textbf{Model-Only (ms)} &
\color{white}\textbf{End-to-End (ms)} &
\color{white}\textbf{Speedup vs.\ Baseline} \\
\midrule
Original FP32 (Baseline)          &  9.91 &  9.20 & 1.00$\times$ \\
\rowcolor{TableAlt}
PyTorch Dynamic INT8              &  9.76 & 10.27 & 1.02$\times$ \\
ONNX Dynamic INT8                 & 15.88 & 18.51 & 0.62$\times$ \\
\rowcolor{TableAlt}
Weight Pruning 50\%               & 13.44 & 14.15 & 0.74$\times$ \\
Weight Pruning 75\%               & 15.58 & 18.97 & 0.64$\times$ \\
\rowcolor{TableAlt}
Weight Sharing k=16               & 14.08 & 15.42 & 0.70$\times$ \\
Weight Sharing k=32               & 11.25 & 14.49 & 0.88$\times$ \\
Sequence Downscale 128 KB         &  4.84 &  5.37 & 2.05$\times$ \\
\rowcolor{TableAlt}
Sequence Downscale 64 KB          &  2.87 &  3.46 & 3.45$\times$ \\
\rowcolor{GreenCell}
\textbf{Sequence Downscale 32 KB} & \textbf{1.82} &  \textbf{2.14} & \textbf{5.44}$\times$ \\
\bottomrule
\end{tabular}
\end{table}

%% ═══════════════════════════════════════════════════════════════════════════
%%  7. MULTI-DIMENSIONAL TRADE-OFF ANALYSIS
%% ═══════════════════════════════════════════════════════════════════════════
\section{Multi-Dimensional Trade-off Analysis}

\subsection{Test-Split Accuracy Across All Compressed Variants}

\begin{tcolorbox}[resultbox]
\textbf{Test-Split Result:} All ten model variants achieved \textbf{100.0\% accuracy
on the 8-sample test split} (4 benign, 4 UPX-packed).  No compression technique
caused any misclassification on the held-out test set.  However, Section~\ref{sec:overfit}
shows that this result must be interpreted alongside the overfitting analysis.
\end{tcolorbox}

\subsection{Memory vs.\ Accuracy Trade-off}

Weight Sharing (k=16) achieves the most dramatic storage compression at
\textbf{521.3 KB -- an 87.5\% reduction} -- while maintaining perfect
classification.  This is because MalConv's byte embedding layer has 257
tokens $\times$ 8 dimensions = 2,056 parameters whose values are highly
correlated (similar byte values have similar byte embeddings).  K-Means with
only 16 centroids captures essentially all the discriminative information.

For reference: a 521 KB model fits comfortably in the 256 KB Flash of an
STM32F4 microcontroller, enabling future deployment on ultra-constrained
embedded targets.

\subsection{Computation vs.\ Latency Trade-off}

Sequence Downscaling demonstrates perfect MAC-to-latency proportionality:

\begin{table}[h!]
\centering
\caption{MAC Reduction vs.\ Latency Speedup (Sequence Downscaling)}
\renewcommand{\arraystretch}{1.3}
\begin{tabular}{C{2.5cm} R{2.8cm} R{2.2cm} R{2.4cm} R{2.4cm}}
\toprule
\rowcolor{TableHead}
\color{white}\textbf{Prefix Size} &
\color{white}\textbf{Active MACs} &
\color{white}\textbf{MAC Reduction} &
\color{white}\textbf{Latency (ms)} &
\color{white}\textbf{Speedup} \\
\midrule
256 KB (baseline) & 536,887,424 & -- & 9.91 ms & 1.00$\times$ \\
\rowcolor{TableAlt}
128 KB & 268,451,968 & 50.0\% & 4.84 ms & 2.05$\times$ \\
64 KB & 134,234,240 & 75.0\% & 2.87 ms & 3.45$\times$ \\
\rowcolor{GreenCell}
\textbf{32 KB}  & \textbf{67,125,376}  & \textbf{87.5\%} & \textbf{1.82 ms} & \textbf{5.44}$\times$ \\
\bottomrule
\end{tabular}
\end{table}

The near-linear relationship between prefix length reduction and latency
reduction confirms that the convolution dominates inference time, and no
other overhead (e.g.\ embedding lookup or pooling) is significant.

\subsection{Pruning -- Computation Without Latency Gains}

A counterintuitive result: pruning at 50\% and 75\% sparsity \textit{increases}
wall-clock latency slightly (13.44 ms and 15.58 ms vs.\ 9.91 ms baseline).
This is because PyTorch's dense convolution operators do not exploit weight
sparsity -- zero-multiplications still consume clock cycles in the dense
arithmetic pipeline.  Latency gains from pruning only materialise on
\textit{sparse-aware hardware} (e.g.\ NVIDIA Ampere sparse cores, or
ARM ML extensions with sparsity support).  The benefit of pruning in this
context is reduced \textit{theoretical} arithmetic complexity (active MACs),
which translates to lower energy consumption and thermal load on embedded
devices, rather than lower wall-clock time on dense CPUs.

%% ─────────────────────────────────────────────────────────────────────────
\subsection{Overfitting Analysis and Generalisation Assessment}
\label{sec:overfit}
%% ─────────────────────────────────────────────────────────────────────────

A critical question raised during evaluation is whether the model's high test
accuracy is a product of genuine generalisation or simply memorisation of the
small training set.  To answer this, the model was re-trained from a random
initialisation (using identical architecture and hyper-parameters) for 30 epochs,
with \textbf{train loss, validation loss, and test loss recorded at every epoch}.
The saved checkpoint was also evaluated separately on all three splits.

\begin{table}[h!]
\centering
\caption{Saved Checkpoint Evaluated on All Three Data Splits}
\renewcommand{\arraystretch}{1.4}
\begin{tabular}{L{3cm} R{3cm} R{3cm}}
\toprule
\rowcolor{TableHead}
\color{white}\textbf{Data Split} &
\color{white}\textbf{BCE Loss} &
\color{white}\textbf{Accuracy} \\
\midrule
Train (52 samples) & 0.0014 & \textbf{100.0\%} \\
\rowcolor{AmberCell}
Validation (6 samples) & 0.3872 & \textbf{83.3\%} \\
\rowcolor{GreenCell}
Test (8 samples) & 0.0516 & \textbf{100.0\%} \\
\midrule
\multicolumn{2}{l}{Generalisation gap (Val loss $-$ Train loss)} & $+$0.3858 \\
\multicolumn{2}{l}{Generalisation gap (Train acc $-$ Val acc)} & $+$16.67\% \\
\bottomrule
\end{tabular}
\end{table}

\begin{tcolorbox}[warnbox]
\textbf{Honest Assessment -- Overfitting is Present.}\quad
The learning curves reveal classic overfitting behaviour:
\begin{itemize}
  \item \textbf{Training loss} collapses to $\approx$0.0000 by epoch~2 and
        remains there for all 30 epochs.
  \item \textbf{Validation loss} drops briefly to 0.45 at epoch~3 then
        \textit{diverges upward} to 0.80 by epoch~30 -- the canonical
        overfitting signature.
  \item \textbf{Validation accuracy} plateaus at \textbf{83.3\%} (5 of 6
        samples correct) while training accuracy is 100\%.
\end{itemize}
The generalisation gap of 0.39 loss units and 16.7 percentage points is
above the standard 0.15-loss threshold for acceptable generalisation.
\end{tcolorbox}

\noindent
The primary causes of overfitting in this project are:
\begin{enumerate}
  \item \textbf{Very small dataset:} Only 66 total samples (33 benign, 33 packed).
        The validation split contains just 6 samples -- too few to provide a
        statistically reliable signal.
  \item \textbf{Highly separable task:} UPX packing produces an extremely strong
        and consistent byte-level signature.  The model saturates quickly and
        then overfits the remaining loss, pushing training scores toward extreme
        values (benign: $<$0.03, packed: $>$0.88).
  \item \textbf{No regularisation:} The baseline training used no dropout, no
        weight decay, and no early stopping.
\end{enumerate}

\noindent
\textbf{Why test accuracy remains 100\%:} The 4 test packed files produce
scores in the range 0.87--0.94; the 4 benign files score 0.0000--0.0014.
Even under significant overfitting, the decision margin is so large that no
sample crosses the 0.5 boundary.  The 100\% test result is \textit{real} for
this specific dataset, but should \textit{not} be extrapolated to unseen
malware families without a larger, more diverse evaluation set.

\subsection{Pareto-Optimal Deployment Choices}

\begin{table}[h!]
\centering
\caption{Deployment Recommendation by Hardware Target}
\renewcommand{\arraystretch}{1.4}
\begin{tabular}{L{3.2cm} L{3cm} L{3.5cm} L{4cm}}
\toprule
\rowcolor{TableHead}
\color{white}\textbf{Target Hardware} &
\color{white}\textbf{Primary Constraint} &
\color{white}\textbf{Recommended Variant} &
\color{white}\textbf{Achieved Metric} \\
\midrule
Raspberry Pi 4 & Latency / throughput & Sequence Downscale 32 KB & 1.82 ms, 100\% acc. \\
\rowcolor{TableAlt}
Raspberry Pi Zero & RAM \& compute & Weight Sharing k=16 + 32 KB downscale & 521 KB, \textasciitilde 1.8 ms \\
STM32F4 (MCU) & Flash \textless 1 MB & Weight Sharing k=16 & 521 KB, 100\% acc. \\
\rowcolor{TableAlt}
Edge Server & Throughput & ONNX Quant + batch inference & INT8 accel., lowest energy \\
\bottomrule
\end{tabular}
\end{table}

%% ═══════════════════════════════════════════════════════════════════════════
%%  8. CONCLUSION
%% ═══════════════════════════════════════════════════════════════════════════
\section{Conclusion and Future Work}

\subsection{Summary of Findings}

\begin{tcolorbox}[resultbox]
\begin{itemize}
  \item The MalConv gated CNN successfully classifies UPX-packed vs.\ benign
        executables using raw byte sequences, achieving \textbf{100\% accuracy
        on the 8-sample test split} without any hand-crafted feature engineering.
  \item \textbf{Overfitting is observed} on the small 66-sample dataset: training
        loss $\rightarrow$ 0, validation loss diverges to 0.80, and validation
        accuracy plateaus at \textbf{83.3\%} -- a 16.7 point gap vs.\ train.
  \item Despite overfitting, \textbf{all four compression techniques preserve
        100\% test-split accuracy}, confirming the model's weight redundancy and
        the large decision-boundary margin of this particular classification task.
  \item \textbf{Best memory compression:} K-Means Weight Sharing (k=16) --
        87.5\% storage reduction to 521 KB with zero test-accuracy loss.
  \item \textbf{Best latency reduction:} Sequence Downscaling to 32 KB --
        5.44$\times$ speedup to 1.82 ms per file inference.
  \item \textbf{Best compute reduction:} Sequence Downscaling to 32 KB --
        87.5\% fewer active MACs (537M $\rightarrow$ 67M).
\end{itemize}
\end{tcolorbox}

\subsection{Project Deliverables}

\begin{table}[h!]
\centering
\caption{Project File Deliverables}
\renewcommand{\arraystretch}{1.3}
\begin{tabular}{L{5.5cm} L{8cm}}
\toprule
\rowcolor{TableHead}
\color{white}\textbf{File} & \color{white}\textbf{Description} \\
\midrule
\texttt{models/malconv/model.py} & MalConv architecture definition \\
\rowcolor{TableAlt}
\texttt{models/malconv/train.py} & Training pipeline with group-split \\
\texttt{models/malconv/export\_malconv\_onnx.py} & ONNX export utility \\
\rowcolor{TableAlt}
\texttt{summary/usb\_reports/optimize\_and\_benchmark.py} & Full optimization \& benchmarking script \\
\texttt{outputs/optimization\_results.json} & Raw benchmark data (JSON) \\
\rowcolor{TableAlt}
\texttt{summary/usb\_reports/tradeoff\_comparison.png} & Static Matplotlib trade-off chart \\
\texttt{summary/usb\_reports/dashboard.html} & Interactive Chart.js web dashboard \\
\rowcolor{TableAlt}
\texttt{outputs/models/group\_split\_model.pth} & Best FP32 baseline checkpoint \\
\texttt{outputs/models/quantized\_dynamic.pth} & PyTorch INT8 checkpoint \\
\rowcolor{TableAlt}
\texttt{outputs/models/shared\_16.pth} & Weight-sharing k=16 checkpoint \\
\texttt{outputs/models/pruned\_50.pth} & 50\% pruned checkpoint \\
\rowcolor{TableAlt}
\texttt{outputs/models/pruned\_75.pth} & 75\% pruned checkpoint \\
\bottomrule
\end{tabular}
\end{table}

\subsection{Future Directions}
\begin{itemize}
  \item \textbf{Address Overfitting with Regularisation:} Add dropout (e.g.\
        $p{=}0.3$ after the gated pooling layer) and L2 weight decay to reduce
        the train-vs-validation gap from 16.7\% toward $<$5\%.
  \item \textbf{Early Stopping:} Monitor validation loss and halt training when
        it begins to diverge (here, that would be epoch~3), preventing the
        model from memorising training-set noise.
  \item \textbf{Larger, Diverse Dataset:} Retrain on EMBER-2018 or SOREL-20M
        with real malware samples to obtain statistically reliable accuracy
        estimates and reduce dependence on any single obfuscation style.
  \item \textbf{Quantization-Aware Training (QAT):} Simulate INT8 arithmetic
        during fine-tuning so quantised weights remain in the optimal range.
  \item \textbf{Structured Pruning:} Prune entire convolutional channels to
        achieve wall-clock latency gains on dense CPUs (unstructured zeros
        do not reduce latency on ARM Cortex without sparse-aware backends).
  \item \textbf{TFLite / CMSIS-NN Export:} Port the 32 KB downscaled + k=16
        shared model to TensorFlow Lite Micro for STM32 Cortex-M4 MCU deployment.
  \item \textbf{YARA + Neural Hybrid:} Combine YARA rule matching (fast, zero-cost
        inference) as a pre-filter with neural scoring, reducing AI inference
        calls by an estimated 80\% on typical USB drives.
\end{itemize}

\end{document}
